
The development of 'foundation models' has transformed natural language processing (NLP) and computer vision. The prevailing training paradigm now employs a two-stage process: pretraining on large unlabeled datasets to learn generalizable representations, followed by fine-tuning on smaller, labeled datasets for specialized tasks. The success of large-scale pretraining has motivated the development of analogous foundation models in atomistic sciences, including chemistry, structural biology, and materials science \cite{AtomisticfFoundationModels_yuan2025, AI4sci_review_zhang2025artificial}. These efforts are driven by their potential for accelerated atomistic simulation, direct structure-to-property prediction, and inverse design—capabilities that could significantly accelerate the discovery of new drugs and materials.

In chemistry and materials science, the most notable progress toward this vision has come from researchers developing machine-learned interatomic potentials (MLIPs) to accelerate atomistic simulation. Pretrained MLIPs - such as UMA\cite{UMA_wood2025}, MACE\cite{MACE_batatia2025} and Orb\cite{Orbv3_rhodes2025} - have achieved remarkable accuracy in DFT force-energy predictions by supervised pretraining on large, diverse atomistic datasets. These models have enabled near-DFT quality simulations and structural relaxations at speeds and scales far exceeding those of traditional DFT. While often termed 'foundation models' (or more appropriately 'foundation potentials'), they differ from their counterparts in NLP and CV in one key regard: current foundation potentials are pretrained exclusively via supervised learning on precomputed forces and energies. However, scaling supervised pretraining with DFT labels is computationally expensive; generating forces and energies for the 100M geometries in the OMol25 dataset required six billion CPU core-hours \cite{OMOL25_levine2025open}.

In other fields, foundation models typically rely on self-supervised learning (SSL) objectives that leverage vast, unlabeled datasets. However, unlike images and text, which are often collected without labels, atomistic datasets are typically generated via computational methods (such as DFT) that inherently produce force and energy labels as part of the data generation process. Although force and energy values vary depending on the DFT functional used, prior work has overcome this by using multiple prediction heads and demonstrating clear benefits to supervised pretraining on large, diverse datasets containing labels derived from multiple functionals \cite{UMA_wood2025, JMP_Shoghi2023}. Consequently, supervised pretraining with DFT labels remains the dominant paradigm, outperforming existing SSL methods on simulation and property prediction benchmarks \cite{JMP_Shoghi2023, NoisyNodes_Zaidi2023PVD}. However, the ratio of 'ground-truth' DFT labeled atomistic data is likely to change as more atomistic geometries are produced using MLIP-driven simulation or relaxation, and generative models \cite{boltz1_2024, boltz2_2025, Sair_dataset, UMA_wood2025, mattergen_zeni2023, MACE_batatia2025, GeomRep_li2024}. Nevertheless, adapting SSL methods from other fields has remained a challenge due to the unique attributes of 3D atomistic data. Atomistic data has no fixed size (samples range from 3 to over 3000 atoms), samples may exhibit periodicity in 3D space, and learned representations must be sensitive to small changes in atomic position to achieve accurate predictions in geometrically sensitive tasks (structural relaxation, simulation, binding affinity predictions, etc). 

Previous work in computer vision suggests that self-supervised pretraining can yield more interpretable and generalizable representations than supervised learning \cite{Mocov3_chen2021, Dino2_oquab2024, DINO3_siméoni2025dinov3}, achieving performance competitive with supervised models even on identical datasets \cite{MAE_he2021, ReLicv2_tomasev2022}. This is because supervised learning only requires a model to represent features of the training data that are relevant to a specific task, becoming invariant to uncorrelated features. In contrast, generative SSL encourages a model to represent the full feature distribution within a training dataset, often producing internal representations with superior transferability. In this work, we demonstrate that this principle holds for 3D atomistic data as well.


% Crucially, atomistic data pose unique challenges when adapting SSL methods from other fields.

% In other fields, foundation models typically rely on self-supervised learning (SSL) objectives that leverage vast, unlabeled datasets. However, unlike images and text, 
% atomistic data poses challenges for SSL. Atomistic data has no fixed size (samples range from 3 to over 3000 atoms), samples may exhibit periodicity in 3D space, and learned representations must be sensitive to small changes in atomic position to achieve accurate predictions in geometrically sensitive tasks (structural relaxation, simulation, binding affinity predictions, etc). 
% Further, compared to language or vision data, a large fraction of currently available atomistic geometries have corresponding DFT-calculated force and energy labels. 
% Unlike images or textual data, which are often collected without labels, atomistic datasets are typically generated via computational methods (such as DFT) that inherently produce force and energy labels as part of the data generation process.
% Although force and energy values vary depending on the DFT functional used, prior work has overcome this by using multiple prediction heads, demonstrating clear benefits to supervised pretraining on large datasets containing labels derived from multiple functionals \cite{UMA_wood2025, JMP_Shoghi2023}. As a result, supervised pretraining with DFT labels remains the dominant approach, currently yielding superior performance for both simulation and property prediction. However, the ratio of DFT labeled atomistic data is likely to change as more atomistic geometries are produced using MLIP-driven simulation or relaxation, and generative models \cite{boltz1_2024, boltz2_2025, Sair_dataset, UMA_wood2025, mattergen_zeni2023, MACE_batatia2025, GeomRep_li2024}.